\section{Future Directions}
\label{sec:gr-future}

Several directions follow from the limitations above and from open problems
identified in the literature:

\begin{itemize}
    \item \textbf{Cheaper, more robust construction.} Reducing the cost and error
    of turning text into graphs---by improving LLM-based extraction and
    linking---remains central to making Graph RAG practical at
    scale~\cite{pan2024unifying, ji2021kgsurvey}.

    \item \textbf{Deeper LLM--KG integration.} Roadmaps for unifying language
    models and knowledge graphs point toward systems where the two are combined
    throughout inference rather than merely at retrieval time, so that generation
    and structured reasoning reinforce each other~\cite{pan2024unifying}.

    \item \textbf{Efficient and incremental indexing.} Making graph indexing
    faster and updatable, as begun by lightweight and incremental
    approaches, is needed for deployment in changing data
    environments~\cite{guo2024lightrag}.

    \item \textbf{Standardized evaluation.} Consolidating reference-free metrics
    for both retrieval quality and answer faithfulness into shared benchmarks
    would let Graph RAG variants be compared reliably~\cite{es2024ragas,
    min2023factscore}.

    \item \textbf{Stronger factual guarantees.} Given the persistence of
    hallucination, using the ontology for verification and provenance is a
    promising route to answers that are not only grounded but
    checkable~\cite{ji2023hallucination, pan2024unifying}.
\end{itemize}
